\documentclass{pgnotes}

\title{Graph databases}

\begin{document}

\maketitle

\section*{Recommended reading}

\begin{enumerate}
\item \href{https://en.wikipedia.org/wiki/Graph_database}{Wikipedia article on graph databases}
\item \href{https://neo4j.com/docs/getting-started/current/}{Neo4j getting started guide}
\end{enumerate}

\section{Graph databases}

Graph databases are a NoSQL database that focus on entities and the relationships between them.
We will look at Neo4j as an example. 

\subsection{Conceptual model}

\begin{description}
\item[Nodes] correspond to entities / rows.
  A node may have:
  \begin{description}
  \item[Properties] for storing information (e.g. age, name)
  \item[Label(s)] for grouping it into a set (e.g. Person, Group)
  \end{description}
\item[Edges / relationships] connect two nodes. Contrasts to a relational (or other noSQL DB) in that relationships are explicitly stored. 
  \begin{itemize}
  \item May have a direction:
    \begin{description}
    \item[Undirected graph] has no directionality implied between the two nodes. (e.g. Friends)
    \item[Directed graph] implies a directionality between the two nodes. (e.g. Manages)
    \end{description}
  \item Edges may have properties independent of the nodes connecting them (e.g. Friend since)
  \end{itemize}
\end{description}

\subsection{Connectivity}

Neo4J supports two access methods: HTTP and Bolt.

\subsection{Query language}

Neo4j uses its own Cypher query language.
SQL is to Postgresql as Cypher is to Neo4j.



\end{document}